\documentclass[main.tex]{subfiles}
\begin{document}\newpage
\setdoublesep{0.35700 em}  % 'Bond Spacing'
\setatomsep{1.78500 em}    % 'Fixed Length'
\setbondoffset{0.18265 em} % 'Margin Width'
\newcommand{\bondwidth}{0.06642 em} % 'Line Width'
\setbondstyle{line width = \bondwidth}

 




%%%%%%%%%%%%HEADING
\begin{multicols}{2}
\begin{tcolorbox}[enhanced jigsaw,breakable,size=title,
colback=mybrown!05,colframe=black,fonttitle=\bfseries,
title=STUDENT INFO,pad at break=1mm, break at=15cm/0pt ]
\vspace{0.2cm}
\noindent Name: \rule{5cm}{0.4pt}Date:\rule{1cm}{0.4pt}\\
Pre-lab Done: \tikzcheckmark[scale=2,black]{no mark}\quad
\end{tcolorbox}
\end{multicols}
\hfill
\vspace{0.2cm}
\begin{center}
{\large \bfseries 
Pre-lab Questions 
\par
\Huge
Nuclear Chemistry
\\[5pt] \par}
\vspace{0.2cm}
\end{center}
\par
\noindent
\uline{  \hfill \normalsize \hfill       }
%%%%%%%%%%%%HEADING

\begin{enumerate}
% PELAB 1
\item Classify the following decays as $\alpha$, $\beta$, $\gamma$ or positron emission:
\begin{enumerate}[label=(\alph*)]
\item \centerline{\ce{^235_92U -> ^231_90Th + ^4_2He} }
\item \centerline{\ce{^235_92U -> ^235_92Np + ^0_{0}\gamma} }
\item \centerline{\ce{^14_6C -> ^14_7N + ^0_{-1}e} }
\item \centerline{\ce{^222_88Ra -> ^4_2He + ^226_{88}Ra} }

\end{enumerate}

\item Iodine-131 is used to treat certain types of thyroid cancer and some rarer types of cancer. Given that its half-life is 8days, calculate the amount of iodine of a 4g sample that remains after 10 days.
\vspace{3cm}
\item Strontium-89  is used to treat some types of secondary bone cancer. Given that its half-life is 50days, calculate the time needed to reduce a 4g sample into 2g.
\vspace{3cm}
\item Iron-59 is used in studies of iron metabolism in the spleen. A Iron-59 sample has an activity of 20cpm and 10cpm 46 days after the first measurement. Calculate the half life of the isotope.
\vspace{3cm}





\end{enumerate}


\clearpage\mbox{}\clearpage



%%%%%%%%%%%%HEADING
\begin{multicols}{2}
\begin{tcolorbox}[enhanced jigsaw,breakable,size=title,
colback=mybrown!05,colframe=black,fonttitle=\bfseries,
title=STUDENT INFO,pad at break=1mm, break at=15cm/0pt ]
\vspace{0.2cm}
\noindent Name: \rule{5cm}{0.4pt}Date:\rule{1cm}{0.4pt}\\
Pre-lab Done: \tikzcheckmark[scale=2,black]{no mark}\quad
\end{tcolorbox}
\end{multicols}
\hfill
\vspace{0.2cm}
\begin{center}
{\large \bfseries 
Experiment
\par
\Huge
Nuclear Chemistry
\\[5pt] \par}
\vspace{0.2cm}
\end{center}
\par
\noindent
\uline{  \hfill \normalsize \hfill       }
%%%%%%%%%%%%HEADING

\vspace{0.2cm}{\large \bfseries 0. Instructions for the use of ST360 radiation counter}
\begin{steps}
    \newstep[] Turn on the radiation counter (red button on the back).
        \newstep[]  Use display select to select \emph{time}. Set the time to 60 seconds by pressing the display Up or Down.
            \newstep[]  Use display select until the light cursor is next to \emph{high voltage}. Set the voltage to 900V  by pressing the display Up or Down.
                \newstep[]  Use display select until the light cursor is next to counts.
                \newstep[]  Place a radioactive chip on the chip support and measure. In case of using a shield, place it between the counter and the chip.
                \newstep[]  Press Count and wait until the stop button sign lights up, as the machine stops automatically when the 60s is up. 
                \newstep[]  Right down the unit of the reading as count per minute (cpm). 
    \end{steps}



\vspace{0.2cm}{\large \bfseries 1. Background radiation}
The air has a certain radioactivity called background radioactivity. This is a very small activity but still affect the radioactive measurements and hence it should be taken into account. In this experiment you will measure the background radiation by means of a Geiger counter. You will have to repeat the measurement several times and average the radiation measured in order to obtain a reliable number.
\begin{steps}
 \newstep[] Do not use any of the chips and make sure they are in the secured protecting box.
    \newstep[] Start the Geiger counter. Set up the measurement time to 60 seconds and the measuring voltage according to your professor's instructions. Mind to select a voltage of 900V for all measurements (Press Display/High Voltage/Up/Down until you reach 900V). Press measure (press Display until the light cursor is next to count; then press Count until the stop button lights up.) and write down the background radioactivity in counts per minute in the table below.
  \newstep[] Repeat the measurement two more times and calculate the average by adding the three measurements and dividing by three. Make sure the measurements are consistent with each other.
\end{steps}

\begin{center}\resizebox{18cm}{!} {\begin{tabular}{ |p{4cm}|p{4cm}|p{4cm}|p{4cm}|  }
\hline
    Measurement 1 (cpm) &   Measurement  2 (cpm) &   Measurement  3 (cpm) &    Average Radiation (cpm)        \\
\hline
   \vspace{0cm}\vspace{.5cm} &     &     &       \\
\hline
\end{tabular}}\end{center}

\vspace{0.2cm}{\large \bfseries 2. Radioactive chips}
In this section you will calculate the radioactivity of a set of different radioactive chops. You will still use the Geiger counter and after measuring the number of counts per minutes you will have to subtract the background radiation to your measurement.

\begin{steps}
    \newstep[] Select three of the radioactive chips.
        \newstep[]  Place one of the chips 5cm away from the counter by means of the plastic stand. 
        \newstep[]  Start the Geiger counter. Set up the measurement time to 60 seconds and the measuring voltage according to your professor's instructions. Press measure and write down the background radioactivity in counts per minute in the table below.
    \newstep[] Repeat the measurement for the other two chips.
    \newstep[] Repeat the measurement for other materials such as tea, instant coffee, potassium chloride or dry seaweed. 
      \newstep[] Now subtract the background radiation measured in the previous section to teach of the measurements.
      
\end{steps}

\begin{center}\begin{tabular}{ |p{3cm}|p{5cm}|p{5cm}|  }
\hline
    Isotope name &   Activity (cpm) &   Activity - Background (cpm)         \\
\hline
   \vspace{0cm}\vspace{.5cm} &     &            \\
\hline
   \vspace{0cm}\vspace{.5cm} &     &            \\
\hline
   \vspace{0cm}\vspace{.5cm} &     &            \\
\hline
   \vspace{0cm}\vspace{.5cm} &     &            \\
\hline
   \vspace{0cm}\vspace{.5cm} &     &            \\
\hline
   \vspace{0cm}\vspace{.5cm} &     &            \\
\hline
   \vspace{0cm}\vspace{.5cm} &     &            \\
\hline
\end{tabular}\end{center}
 


 
 \vspace{0.2cm}{\large \bfseries 3. Radioactivity protection}
In this section you will only use one of the chips from the previous experiment. For this one chip you will use different barriers to shield radiation and estimate the shielding impact.
\begin{steps}
    \newstep[] Select radioactive chip from the previous experiment that gave you the highest counts per minute.
        \newstep[]  Place the chips at a 10cm distance from the counter by means of the plastic stand. 
        \newstep[]  Select three of the shielding and place one of these in between the counter and the chip.
    \newstep[] Measure the number of counts per minute and subtract the background radiation.
   \newstep[] Write down the measurement in the table below. Compute the activity taking into account the background radiation. These results can potentially be a negative value.
       \newstep[] Repeat the procedure for the other two shielding.      
\end{steps}

\begin{center}Isotope name=\rule{5cm}{0.4pt}\end{center}
\begin{center}\begin{tabular}{ |p{3cm}|p{5cm}|p{5cm}|  }
\hline
    Shielding name &   Activity (cpm) &   Activity - Background (cpm)         \\
\hline
   \vspace{0cm}\vspace{.5cm} &     &            \\
\hline
   \vspace{0cm}\vspace{.5cm} &     &            \\
\hline
   \vspace{0cm}\vspace{.5cm} &     &            \\
\hline
\end{tabular}\end{center}



\vspace{0.2cm}{\large \bfseries 5. Effect of distance on radioactivity }
In this section you will only use one of the chips from the previous experiment, again the stronger one. You will place this chip at different distances from the counter and measure the impact of distance on radioactivity.
\begin{steps}
        \newstep[]  Place the chips at five different distances from the counter by means of the plastic stand. The distances are indicated in the table below.
    \newstep[] Measure the number of counts per minute for each distance and subtract the background radiation.
      \newstep[] Write down the measurement in the table below. 
            \newstep[] Plot activity without the background (right column, vertical axis) vs distance (horizontal axis) connecting the points with a line in the graph below. 

\end{steps}

\begin{center}Isotope name=\rule{5cm}{0.4pt}\end{center}
\begin{center}\begin{tabular}{ |p{3cm}|p{5cm}|p{5cm}|  }
\hline
    Distance (cm) &   Activity (cpm) &   Activity - Background (cpm)         \\
\hline
   \vspace{0cm}\vspace{.5cm}1cm &     &            \\
\hline
   \vspace{0cm}\vspace{.5cm}2cm &     &            \\
\hline
   \vspace{0cm}\vspace{.5cm}3cm &     &            \\
\hline
   \vspace{0cm}\vspace{.5cm}4cm &     &            \\
\hline
   \vspace{0cm}\vspace{.5cm}5cm &     &            \\
\hline
\end{tabular}\end{center}



 


\begin{center}
 \begin{tikzpicture}
 \draw[step=1mm,help lines] (0,0) grid (120mm,120mm);
 \draw[step=10mm] (0,0) grid (120mm,120mm);
\end{tikzpicture}\\
\vspace{0.2cm}
Isotope Name=\rule{5cm}{0.4pt}
\end{center}

%\clearpage\mbox{}\clearpage

%%%%%%%%%%%%HEADING
\newpage\begin{multicols}{2}
\begin{tcolorbox}[enhanced jigsaw,breakable,size=title,
colback=mybrown!05,colframe=black,fonttitle=\bfseries,
title=STUDENT INFO,pad at break=1mm, break at=15cm/0pt ]
\vspace{0.2cm}
\noindent Name: \rule{5cm}{0.4pt}Date:\rule{1cm}{0.4pt}\\
Pre-lab Done: \tikzcheckmark[scale=2,black]{no mark}\quad
\end{tcolorbox}
\end{multicols}
\hfill
\vspace{0.2cm}
\begin{center}
{\large \bfseries 
Post-lab Questions 
\par
\Huge
Nuclear Chemistry
\\[5pt] \par}
\vspace{0.2cm}
\end{center}
\par
\noindent
\uline{  \hfill \normalsize \hfill       }
%%%%%%%%%%%%HEADING

\begin{enumerate}
\item From the radioactive chips you studied indicate the nature of the radiation (written on the chip) produced by the strongest chip.
\vspace{2.5cm}

\item From the common materials, you tested for radiation (tea, seaweed, coffee, and KCl), which one gave you the highest radioactive measurement?
\vspace{2.5cm}

\item From the different shielding you studied indicate the nature of the one that protected the most from radiation.
\vspace{2.5cm}

\item From the graph you made estimate the number of counts per minute at 3.5cm from the counter.
\vspace{2.5cm}



\end{enumerate}

\clearpage\mbox{}\clearpage


 


\end{document}
